% Top-level report for the Campbell-Shiller replication pipeline.
% This file lives next to the other samples so students can see the
% full document, but it only compiles inside topic.04/sample4_empirical/
% latex/ where the pipeline (one dir up) produces ../outputs/yields.png,
% ../outputs/summary.tex, and ../outputs/cs_table.tex.

\documentclass{article}

\usepackage{amsmath, amssymb}
\usepackage{graphicx}
\usepackage{booktabs}
\usepackage{siunitx}
\usepackage{threeparttable}
\usepackage[margin=1in]{geometry}

\title{The Campbell-Shiller Puzzle: \\
       A Replication with GSW Fitted Zero-Coupon Yields}
\author{FBE 633 --- Topic 4 Example}
\date{\today}

\begin{document}
\maketitle

\section{Introduction}

Following Dai and Singleton (2002, Table 1), we estimate the Campbell-Shiller
long-rate regression
\begin{equation}
  R^{n-1}_{t+1} - R^n_t = \alpha_n + \phi_n \cdot \frac{R^n_t - r_t}{n - 1} + \varepsilon_{t+1},
  \label{eq:cs}
\end{equation}
where $R^n_t$ is the $n$-year zero-coupon yield and $r_t \equiv R^1_t$ is the short rate.
Under the expectations hypothesis, $\phi_n = 1$ for all $n$.

We use Gurkaynak-Sack-Wright fitted zero-coupon Treasury yields, monthly,
1985--present, with $n \in \{2, 3, \ldots, 10\}$ years.

\section{Yield Data}

\begin{figure}[!h]
  \centering
  \includegraphics[width=0.95\linewidth]{../outputs/yields.png}
  \caption{GSW fitted zero-coupon Treasury yields, 1985--present.}
  \label{fig:yields}
\end{figure}

\input{../outputs/summary.tex}

\section{Campbell-Shiller Regression}

\input{../outputs/cs_table.tex}

The classical pattern reappears: $\phi_n$ is significantly below unity at
every maturity, and becomes more negative with $n$. The expectations
hypothesis is rejected, and the rejection deepens with maturity.

\end{document}
